%!TEX root = forallxyyc.tex
\part{Metateoria}
\label{ch.normalform}
\addtocontents{toc}{\protect\mbox{}\protect\hrulefill\par}

\chapter{Formas normais}\label{c:NormalForms}


\section{Forma normal disjuntiva}\label{s:DNFDefined}

Às vezes é útil considerar sentenças de uma forma particularmente simples. Por exemplo, podemos considerar sentenças nas quais $\enot$ se aplica apenas a sentenças atômicas, ou combinações de sentenças atômicas e sentenças atômicas negadas usando somente $\eand$. Uma forma relativamente geral, mas ainda simples, é aquela em que uma sentença é uma disjunção de conjunções de sentenças atômicas ou sentenças atômicas negadas. Para construir uma sentença assim, começamos com sentenças atômicas, depois (talvez) acrescentamos negações, em seguida (talvez) combinamos usando~$\eand$ e, por fim (talvez), combinamos usando~$\eor$.

Diremos que uma sentença está na \define{forma normal disjuntiva} \emph{\ifeff} satisfaz todas as condições a seguir:
	\begin{compactlist}
		\item[(\textsc{dnf1})] Não ocorrem na sentença conectivos além de negações, conjunções e disjunções;
		\item[(\textsc{dnf2})] Toda ocorrência de negação tem escopo mínimo (isto é, qualquer ocorrência de $\enot$ é imediatamente seguida por uma sentença atômica);
		\item[(\textsc{dnf3})] Nenhuma disjunção ocorre no escopo de uma conjunção.
	\end{compactlist}
\newglossaryentry{disjunctive normal form}{
  name = forma normal disjuntiva (FND),
  text = forma normal disjuntiva,
  description = {Uma sentença que é uma disjunção de conjunções de sentenças atômicas ou sentenças atômicas negadas}
}
Eis algumas sentenças na forma normal disjuntiva:
\begin{align*}
  & A \\
  & (A \eand \lnot B \eand C)\\
  & (A \eand B) \eor (A \eand \enot B)\\
  & (A \eand B) \eor (A \eand  B \eand C \eand \enot D \eand \enot E)\\
  & A \eor (C \eand \enot P_{234} \eand P_{233} \eand Q) \eor \enot B
\end{align*}
Observe que aqui violamos um dos preceitos deste livro e \emph{temporariamente} permitimos as convenções relaxadas de uso de parênteses, que admitem que conjunções e disjunções tenham comprimento arbitrário. Essas convenções facilitam ver quando uma sentença está na forma normal disjuntiva. Continuaremos a usá-las sem mais comentários.

Para ilustrar melhor a ideia de forma normal disjuntiva, introduziremos mais uma notação. Escrevemos o símbolo $\pm\metav{A}$ para indicar que $\metav{A}$ é uma sentença atômica que pode ou não ser precedida por uma ocorrência de negação. Então, uma sentença na forma normal disjuntiva tem o seguinte formato:
	$$(\pm \metav{A}_1 \land \ldots \land \pm \metav{A}_i) \lor (\pm \metav{A}_{i+1} \land \ldots \land \pm\metav{A}_j) \lor \ldots \lor (\pm\metav{A}_{m+1} \land \ldots \land \pm \metav{A}_n)$$
Agora sabemos o que significa uma sentença estar na forma normal disjuntiva. O resultado que pretendemos alcançar é:
	\factoidbox{\label{thm:dnf}\textbf{Teorema da forma normal disjuntiva.} Para toda sentença, existe uma sentença equivalente na forma normal disjuntiva.
	}
Daqui em diante, abreviaremos «forma normal disjuntiva» por «FND».


\section{Demonstração do teorema da FND por meio de tabelas-verdade}
\label{s:DNFTruthTable}

Nossa primeira demonstração do teorema da FND emprega tabelas-verdade. Primeiro ilustraremos a técnica para encontrar uma sentença equivalente na FND e, em seguida, transformaremos essa ilustração em uma demonstração rigorosa.

Suponhamos que temos uma sentença, $\metav{S}$, que contém três sentenças atômicas: $A$, $B$ e $C$. A primeira coisa a fazer é preencher uma tabela-verdade completa para~$\metav{S}$. Talvez obtenhamos esta:
\begin{center}
\begin{tabular}{c c c | c}
$A$ & $B$ & $C$ & $\metav{S}$\\
\hline
 T & T & T & T \\
 T & T & F & F \\
 T & F & T & T \\
 T & F & F & F \\
 F & T & T & F \\
 F & T & F & F \\
 F & F & T & T \\
 F & F & F & T
\end{tabular}
\end{center}
%Now, consider a sentence, whose only connectives are negations and conjunctions, where no connective occurs within the scope of any negation, e.g.:
%	$$A \eand \enot B \eand C$$
%This sentence is true when, and only when, `$A$' is true, `$B$' is false and `$C$' is true. Similarly, the sentence:
%	$$\enot A \eand \enot B \eand C$$
%this is true when, and only when, `$A$' is false, `$B$' is false and `$C$' is true.
%
%A disjunction is true when, and only when, at least one of the disjuncts is true. So if we write down a disjunction of sentences of the above form, perhaps
%	$$(A \eand \enot B \eand C) \eor (\enot A \eand \enot B \eand C)$$
%then it will be true on exactly \emph{two} lines of the truth table which describes all possible valuations of `$A$', `$B$' and `$C$'.
%
Acontece que $\metav{S}$ é verdadeira em quatro linhas de sua tabela-verdade, a saber, as linhas 1, 3, 7 e~8. Para cada uma dessas linhas, escreveremos uma sentença cujos únicos conectivos são negações e conjunções, e na qual toda negação tem escopo mínimo:
	\begin{compactlist}
		\item $A \eand B \eand C$\hfill que é verdadeira na linha 1 (e somente nela)
		\item $A \eand \enot B \eand C$ \hfill que é verdadeira na linha 3 (e somente nela)
		\item $\enot A \eand \enot B \eand C$ \hfill que é verdadeira na linha 7 (e somente nela)
		\item $\enot A \eand \enot B \eand \enot C$ \hfill que é verdadeira na linha 8 (e somente nela)
	\end{compactlist}
Agora combinamos todas essas conjunções usando~$\eor$, assim:
$$(A \eand B \eand C) \eor (A \eand \enot B \eand C) \eor (\enot A \eand \enot B \eand C) \eor (\enot A \eand \enot B \eand \enot C)$$\label{longDNF}
Isso nos dá uma sentença na FND que é verdadeira exatamente nas linhas em que uma das disjunções é verdadeira, isto é, nas linhas 1, 3, 7 e 8 (e somente nelas). Portanto, essa sentença tem exatamente a mesma tabela-verdade que $\metav{S}$. Temos, assim, uma sentença na FND logicamente equivalente a $\metav{S}$, exatamente o que queríamos!

A estratégia que acabamos de adotar não dependia das particularidades de $\metav{S}$; ela é perfeitamente geral. Consequentemente, podemos usá-la para obter uma demonstração simples do teorema da FND.

Escolha uma sentença arbitrária, $\metav{S}$, e sejam $\metav{A}_1, \ldots, \metav{A}_n$ as sentenças atômicas que ocorrem em $\metav{S}$. Para obter uma sentença na FND logicamente equivalente a $\metav{S}$, consideramos a tabela-verdade de $\metav{S}$. Há dois casos a considerar:
	\begin{numberlist}
		\item \emph{$\metav{S}$ é falsa em todas as linhas de sua tabela-verdade.} Então $\metav{S}$ é uma contradição. Nesse caso, a contradição $(\metav{A}_1 \eand \enot \metav{A}_1)$ está na FND e é logicamente equivalente a~$\metav{S}$.

		\item \emph{$\metav{S}$ é verdadeira em pelo menos uma linha de sua tabela-verdade.}
		Para cada linha $i$ da tabela-verdade, seja $\metav{B}_i$ uma conjunção da forma
		$$(\pm\metav{A}_1 \land \ldots \land \pm\metav{A}_n)$$
		em que as regras a seguir determinam se devemos ou não incluir uma negação diante de cada sentença atômica:
			\begin{itemize}
				\item
				$\metav{A}_m$ é um conjuntivo de $\metav{B}_i$ \emph{\ifeff}
				$\metav{A}_m$ é verdadeira na linha~$i$.
				\item $\enot\metav{A}_m$ é um conjuntivo de $\metav{B}_i$
				\emph{\ifeff} $\metav{A}_m$ é falsa na linha~$i$.
			\end{itemize}
		Dadas essas regras, $\metav{B_i}$ é verdadeira na linha~$i$ (e somente nela) da tabela-verdade que considera todas as valorações possíveis de $\metav{A}_1,\ldots,\metav{A}_n$ (isto é, a tabela-verdade de $\metav{S}$).

		Agora sejam $\metav{ i_1}$, $\metav{ i_2}$, \dots, $\metav{ i_m}$ os números das linhas da tabela-verdade em que $\metav{S}$ é \emph{verdadeira}. Seja então $\metav{D}$ a sentença:
		$$\metav{B}_{i_1} \eor \metav{B}_{i_2} \eor \ldots \eor \metav{B}_{i_m}$$
		Como $\metav{S}$ é verdadeira em pelo menos uma linha de sua tabela-verdade, $\metav{D}$ é de fato bem definida; e, no caso-limite em que $\metav{S}$ é verdadeira em exatamente uma linha de sua tabela-verdade, $\metav{D}$ é simplesmente $\metav{B}_{i_1}$, para algum $i_1$.

		Por construção, $\metav{D}$ está na FND. Além disso, por construção, para cada linha~$i$ da tabela-verdade: $\metav{S}$ é verdadeira na linha $i$ da tabela-verdade \emph{\ifeff} uma das disjunções de $\metav{D}$ (a saber, $\metav{B_i}$) é verdadeira na linha $i$ e somente nela. Logo, $\metav{S}$ e $\metav{D}$ têm a mesma tabela-verdade e, portanto, são logicamente equivalentes.
	\end{numberlist}
Os dois casos são exaustivos e, em qualquer deles, temos uma sentença na FND logicamente equivalente a $\metav{S}$.

Assim, demonstramos o teorema da FND. Antes de prosseguir, porém, devemos assinalar imediatamente que estamos retornando à definição austera de sentença da LVF, segundo a qual podemos supor que toda conjunção tem exatamente dois conjuntivos e toda disjunção tem exatamente dois disjuntivos.


\section{Forma normal conjuntiva}
\label{s:CNF}

Até aqui, neste capítulo, discutimos a forma normal \emph{disjuntiva}. Talvez não seja surpresa saber que também existe a forma normal \emph{conjuntiva} (FNC).

A definição de FNC é exatamente análoga à definição de FND. Assim, uma sentença está na FNC \emph{\ifeff} satisfaz todas as condições a seguir:
	\begin{compactlist}
		\item[(\textsc{cnf1})] Não ocorrem na sentença conectivos além de negações, conjunções e disjunções;
		\item[(\textsc{cnf2})] Toda ocorrência de negação tem escopo mínimo;
		\item[(\textsc{cnf3})] Nenhuma conjunção ocorre no escopo de uma disjunção.
	\end{compactlist}
\newglossaryentry{conjunctive normal form}{
  name = forma normal conjuntiva (FNC),
  text = forma normal conjuntiva,
  description = {Uma sentença que é uma conjunção de disjunções de sentenças atômicas ou sentenças atômicas negadas}
}
Em geral, portanto, uma sentença na FNC tem este aspecto
	$$(\pm \metav{A}_1 \lor \ldots \lor \pm \metav{A}_i) \land (\pm \metav{A}_{i+1} \lor \ldots \lor \pm\metav{A}_j) \land \ldots \land (\pm\metav{A}_{m+1} \lor\ldots \lor \pm \metav{A}_n)$$
onde cada $\metav{A}_k$ é uma sentença atômica.

Podemos agora demonstrar outro teorema da forma normal:
	\factoidbox{\label{thm:cnf}\textbf{Teorema da forma normal conjuntiva.} Para toda sentença, existe uma sentença equivalente na forma normal conjuntiva.}


	Para uma sentença da LVF, $\metav{S}$, começamos escrevendo a tabela-verdade completa de $\metav{S}$.

	Se $\metav{S}$ é \emph{verdadeira} em todas as linhas da tabela-verdade, então $\metav{S}$ e $(\metav{A}_1 \eor \enot \metav{A}_1)$ são logicamente equivalentes.

	Se $\metav{S}$ é \emph{falsa} em pelo menos uma linha da tabela-verdade, então, para cada linha em que $\metav{S}$ é falsa, escreva uma disjunção $(\pm\metav{A}_1 \eor \ldots \eor \pm\metav{A}_n)$ que seja \emph{falsa} nessa linha (e somente nela). Seja $\metav{C}$ a conjunção de todas essas disjunções; por construção, $\metav{C}$ está na FNC e $\metav{S}$ e $\metav{C}$ são logicamente equivalentes.

\practiceproblems
\problempart
\label{pr.DNF}
Considere as seguintes sentenças:
	\begin{compactlist}
		\item $(A \eif \enot B)$
		\item $\enot (A \eiff B)$
		\item $(\enot A \eor \enot (A \eand B))$
		\item $(\enot (A \eif B ) \eand (A \eif C))$
		\item $(\enot (A \eor B) \eiff ((\enot C \eand \enot A) \eif \enot B))$
		\item $((\enot (A \eand \enot B) \eif C) \eand \enot (A \eand D))$
	\end{compactlist}
Para cada sentença, encontre uma sentença equivalente na FND e outra na FNC.

\chapter{Completude funcional}\label{c:FunctionalCompleteness}

Entre nossos conectivos, $\enot$ se aplica a uma única sentença, e os demais combinam exatamente duas sentenças. Também podemos introduzir a ideia de um conectivo de $n$ lugares. Por exemplo, podemos considerar um conectivo de três lugares, $\heartsuit$, e estipular que ele tenha a seguinte tabela-verdade característica:
\begin{center}
\begin{tabular}{c c c | c}
$A$ & $B$ & $C$ & $\heartsuit(A,B,C)$\\
\hline
 T & T & T & F \\
 T & T & F & T \\
 T & F & T & T \\
 T & F & F & F \\
 F & T & T & F \\
 F & T & F & T \\
 F & F & T & F \\
 F & F & F & F
\end{tabular}
\end{center}
Provavelmente esse novo conectivo não corresponderia a nenhuma expressão natural em português (pelo menos não da maneira como $\eand$ corresponde a «e»). Mas surge uma pergunta: se quiséssemos empregar um conectivo com essa tabela-verdade característica, teríamos de acrescentar um conectivo \emph{novo} à LVF? Ou podemos nos virar com os conectivos que \emph{já temos} (como podemos fazer com o conectivo «nem\dots nem», por exemplo)?

Tornemos essa pergunta mais precisa. Diremos que alguns conectivos são \define{conjuntamente funcionalmente completos} \emph{\ifeff} para toda tabela-verdade possível, existe uma sentença que contém somente esses conectivos e tem essa tabela-verdade.

\newglossaryentry{functionally complete}{
  name = completude funcional,
  text = funcionalmente completo,
  description = {Propriedade de uma coleção de conectivos que ocorre \ifeff{} toda tabela-verdade possível é a tabela-verdade de uma sentença que envolve somente esses conectivos}}

A ideia geral é que, munidos de alguns conectivos conjuntamente funcionalmente completos, nenhuma tabela-verdade característica fica fora do nosso alcance. E, de fato, temos sorte.
	\factoidbox{\label{thm:ExpressiveAdequacy}\textbf{Teorema da completude funcional.}
Os conectivos da LVF são conjuntamente funcionalmente completos. De fato, os seguintes pares de conectivos são conjuntamente funcionalmente completos:
\begin{compactlist}
\item\label{expressive:eor} $\enot$ e $\eor$
\item\label{expressive:eand} $\enot$ e $\eand$
\item\label{expressive:eif} $\enot$ e $\eif$
\end{compactlist}}

Dada qualquer tabela-verdade, podemos usar o método da demonstração do teorema da FND (ou do teorema da FNC) por meio de tabelas-verdade, em \cref{c:NormalForms}, para escrever um esquema que tenha a mesma tabela-verdade. Por exemplo, empregando o método das tabelas-verdade para demonstrar o teorema da FND, descobrimos que o esquema a seguir tem a mesma tabela-verdade característica que $\heartsuit(A,B,C)$ acima:
		$$(A \eand B \eand \enot C) \eor (A \eand \enot B \eand C) \eor (\enot A \eand B \eand \enot C)$$
Segue-se que os conectivos da LVF são conjuntamente funcionalmente completos. Agora demonstramos cada um dos resultados auxiliares.

\emph{Resultado auxiliar 1: completude funcional de $\enot$ e $\eor$.} Observe que o esquema que geramos usando o método das tabelas-verdade para demonstrar o teorema da FND conterá somente os conectivos $\enot$, $\eand$ e $\eor$. Portanto, basta mostrar que existe um esquema equivalente que contenha somente $\enot$ e $\eor$. Para isso, basta observar que
		\begin{align*}
		(\metav{A} \eand \metav{B}) & \Leftrightarrow \enot(\enot \metav{A} \eor\enot \metav{B})
		\end{align*}
são logicamente equivalentes.\footnote{Aqui usamos a convenção de que dizer que $\metav{A}$ e $\metav{B}$ são equivalentes pode ser abreviado como $\metav{A} \Leftrightarrow \metav{B}$. Observe que $\Leftrightarrow$ \emph{não} é um conectivo como~$\eiff$, mas um símbolo metalinguístico, assim como~$\entails$; compare a discussão de $\entails$ e $\eiff$ em \cref{entails-v-iff}.}






\emph{Resultado auxiliar 2: completude funcional de $\enot$ e $\eand$.} Exatamente como no resultado auxiliar~1, fazendo uso do fato de que
		\begin{align*}
		(\metav{A} \eor \metav{B}) & \Leftrightarrow \enot(\enot \metav{A} \eand\enot \metav{B})
		\end{align*}
são logicamente equivalentes.

\emph{Resultado auxiliar 3: completude funcional de $\enot$ e $\eif$.} Exatamente como no resultado auxiliar~1, fazendo uso destas equivalências:
		\begin{align*}
		(\metav{A} \eor \metav{B}) &\Leftrightarrow (\enot \metav{A} \eif \metav{B})\\
		(\metav{A} \eand \metav{B}) &\Leftrightarrow \enot(\metav{A} \eif \enot\metav{B})
		\end{align*}
Alternativamente, poderíamos simplesmente usar um dos outros dois resultados auxiliares e invocar (repetidamente) apenas uma dessas duas equivalências.

Em suma, nunca há \emph{necessidade} de acrescentar novos conectivos à LVF. Na verdade, já existe certa redundância entre os conectivos que temos: poderíamos ter nos contentado com apenas dois conectivos, se estivéssemos nos sentindo realmente austeros.

\ifHTMLtarget
\section{Conectivos individualmente funcionalmente completos}
\else
\section[Conectivos funcionalmente completos][Conectivos funcionalmente completos]{Conectivos individualmente funcionalmente completos}
\fi

De fato, alguns conectivos de dois lugares são \emph{individualmente} funcionalmente completos. Esses conectivos normalmente não são incluídos na LVF, pois são bastante trabalhosos de usar. Mas sua existência mostra que, se quiséssemos, poderíamos ter definido uma linguagem funcional-veritativa completa contendo apenas um conectivo primitivo.

O primeiro conectivo desse tipo que consideraremos é $\uparrow$, cuja tabela-verdade característica é a seguinte.
\begin{center}
\begin{tabular}{c c | c}
$\metav{A}$ & $\metav{B}$ & $\metav{A} \mathrel{\uparrow} \metav{B}$\\
\hline
 T & T & F \\
 T & F & T \\
 F & T & T  \\
 F & F & T
\end{tabular}
\end{center}
Isso é frequentemente chamado de «barra de Sheffer», em homenagem a Henry Sheffer, que a usou para mostrar como reduzir o número de conectivos lógicos nos \textit{Principia Mathematica}.\footnote{Sheffer, ``A set of five independent postulates for Boolean algebras, with application to logical constants'', \textit{Transactions of the American Mathematical Society}~14 (1913), pp.~481--488} (Na verdade, Charles Sanders Peirce havia antecipado Sheffer em cerca de 30 anos, mas nunca publicou seus resultados, e o filósofo polonês Edward Stamm publicou o mesmo resultado dois anos antes de Sheffer.)\footnote{Veja o texto de Peirce ``A Boolian algebra with one constant'', datado de c.~1880, em \textit{Collected Papers}, vol.~4, pp.~264--5, e Edward Stamm, ``\foreignlanguage{german}{Beitrag zur Algebra der Logik}'', \foreignlanguage{german}{\textit{Monatshefte für Mathematik und Physik}}~22 (1911), pp.~137--49.} Também é bastante comum chamá-lo de «nand», pois sua tabela-verdade característica é a negação da tabela-verdade de `$\eand$'.
















\factoidbox{\label{prop:upexpressive}$\uparrow$ é funcionalmente completo por si só.}

O teorema da completude funcional nos diz que $\enot$ e $\eor$ são conjuntamente funcionalmente completos. Portanto, basta mostrar que, dado qualquer esquema que contenha somente esses dois conectivos, podemos reescrevê-lo como um esquema equivalente que contenha somente $\uparrow$. Assim como na demonstração dos casos auxiliares do teorema da completude funcional, basta aplicar as seguintes equivalências:
		\begin{align*}
			\enot \metav{A} &\Leftrightarrow (\metav{A} \uparrow \metav{A})\\
			(\metav{A} \eor \metav{B}) & \Leftrightarrow ((\metav{A} \uparrow \metav{A}) \uparrow (\metav{B} \uparrow \metav{B}))
		\end{align*}
ao resultado auxiliar~\ref*{expressive:eor}.

De modo semelhante, podemos considerar o conectivo $\downarrow$:
\begin{center}
\begin{tabular}{c c | c}
$\metav{A}$ & $\metav{B}$ & $\metav{A} \mathrel{\downarrow} \metav{B}$\\
\hline
 T & T & F \\
 T & F & F  \\
 F & T & F  \\
 F & F & T
\end{tabular}
\end{center}
Às vezes ele é chamado de «seta de Peirce» (o próprio Peirce o chamou de «ampheck»). Mais frequentemente, porém, é chamado de «nor», pois sua tabela-verdade característica é a negação de $\eor$, isto é, de «nem\dots{} nem\dots».
	\factoidbox{
		$\downarrow$ é funcionalmente completo por si só. }

Assim como no resultado anterior para $\uparrow$, embora invoquemos as equivalências
		\begin{align*}
			\enot \metav{A} &\Leftrightarrow (\metav{A} \downarrow \metav{A})\\
			(\metav{A} \eand \metav{B}) & \Leftrightarrow ((\metav{A} \downarrow \metav{A}) \downarrow (\metav{B} \downarrow \metav{B}))
		\end{align*}
e o resultado auxiliar~\ref*{expressive:eand}.


\section{Falhas da completude funcional}

Na verdade, os \emph{únicos} conectivos de dois lugares individualmente funcionalmente completos são $\uparrow$ e $\downarrow$. Mas como mostraríamos isso? De modo mais geral, como podemos mostrar que alguns conectivos \emph{não} são conjuntamente funcionalmente completos?

O procedimento óbvio é tentar encontrar alguma tabela-verdade que \emph{não possamos} expressar usando apenas os conectivos dados. Mas há certa arte nisso.

Para tornar isso concreto, consideremos a questão de saber se $\eor$ é funcionalmente completo por si só. Depois de refletir um pouco, deve ficar claro que não é. Em particular, deve ficar claro que nenhum esquema que contenha somente disjunções pode ter a mesma tabela-verdade que a negação, isto é:
				\begin{center}
				\begin{tabular}{c | c}
				$\metav{A}$ & $\enot \metav{A}$\\
				\hline
				 T & F \\
				 F & T
				\end{tabular}
				\end{center}
A razão intuitiva para isso é simples: a primeira linha da tabela-verdade desejada precisa ter o valor Falso; mas a primeira linha de qualquer tabela-verdade de um esquema que contenha \emph{somente} $\eor$ será sempre Verdadeira. O mesmo vale para $\eand$, $\eif$ e $\eiff$.
 	\factoidbox{
		$\eor$, $\eand$, $\eif$ e $\eiff$ não são funcionalmente completos por si sós.}

Na verdade, vale o seguinte:

        \factoidbox{Os \emph{únicos} conectivos de dois lugares funcionalmente completos por si sós são $\uparrow$ e $\downarrow$. }

É claro que isso é mais difícil de demonstrar do que no caso dos conectivos primitivos. Por exemplo, o conectivo «ou exclusivo» não tem um T na primeira linha de sua tabela-verdade característica, e portanto o método usado acima já não basta para mostrar que ele não pode expressar todas as tabelas-verdade. Também é mais difícil mostrar que, por exemplo, $\eiff$ e $\enot$ \emph{juntos} não são funcionalmente completos.

\chapter{Demonstração de equivalências}\label{ch:equivalences}

\section{Substituibilidade de equivalentes}

Lembre-se, de \cref{sec:equivalent}, que $\metav{P}$ e $\metav{Q}$ são equivalentes (na LVF) se, e somente se, para toda valoração, seus valores de verdade coincidem. Vimos muitos exemplos disso e usamos tanto tabelas-verdade quanto demonstrações de dedução natural para estabelecer tais equivalências. Em \cref{c:NormalForms}, chegamos a demonstrar que toda sentença da LVF é equivalente a uma sentença na forma normal conjuntiva e a uma na forma normal disjuntiva. Se $\metav{P}$ e~$\metav{Q}$ são equivalentes, elas sempre têm o mesmo valor de verdade, cada uma implica a outra e, de qualquer uma delas, podemos demonstrar a outra.

Sentenças equivalentes não são, é claro, a mesma sentença: as sentenças $\enot\enot A$ e~$A$ podem sempre ter o mesmo valor de verdade, mas a primeira começa com o símbolo $\enot$, enquanto a segunda não. Mas você pode se perguntar se é sempre verdade que podemos substituir uma das sentenças de um par de sentenças equivalentes pela outra e se os resultados também serão equivalentes. Por exemplo, considere $\enot\enot A \eif B$ e $A \eif B$. A segunda resulta da primeira substituindo $\enot\enot A$ por~$A$. E essas duas sentenças também são equivalentes.

Esse é um fato geral, e não é difícil ver por que ele é verdadeiro. Em qualquer valoração, calculamos o valor de verdade de uma sentença «de dentro para fora». Assim, ao determinar o valor de verdade de $\enot\enot A \eif B$, primeiro calculamos o valor de verdade de $\enot\enot A$, e o valor da sentença completa depende então apenas desse valor (verdadeiro ou falso, conforme o caso) e do restante da sentença (o valor de verdade de~$B$ e a tabela-verdade de~$\eif$). Mas, como $\enot\enot A$ e~$A$ são equivalentes, elas sempre têm o mesmo valor de verdade em uma dada valoração; portanto, substituir $\enot\enot A$ por~$A$ não pode alterar o valor de verdade da sentença completa. O mesmo vale, é claro, para qualquer outra sentença equivalente a $\enot\enot A$, digamos, $A \eand (A \eor A)$.












Para enunciar o resultado em geral, usemos a notação $\metav{R}(\metav{P})$ para indicar uma sentença que contém a sentença $\metav{P}$ como parte. Então $\metav{R}(\metav{Q})$ significa o resultado de substituir a ocorrência de $\metav{P}$ pela sentença~$\metav{Q}$. Por exemplo, se $\metav{P}$ é a letra sentencial $A$, $\metav{Q}$ é a sentença $\enot\enot A$ e $\metav{R}(\metav{P})$ é $A \eif B$, então $\metav{R}(\metav{Q})$ é $\enot\enot A \eif B$.

\factoidbox{Se $\metav{P}$ e $\metav{Q}$ são equivalentes, então também são equivalentes $\metav{R}(\metav{P})$ e $\metav{R}(\metav{Q})$.}

Segue-se desse fato que qualquer sentença da forma $\metav{R}(\metav{P}) \eiff \metav{R}(\metav{Q})$, em que $\metav{P}$ e $\metav{Q}$ são equivalentes, é uma tautologia. No entanto, as demonstrações por dedução natural serão muito diferentes para diferentes $\metav{R}$. (Como exercício, dê demonstrações que mostrem que
\begin{align*}
	& \vdash (\enot\enot P \eif Q) \eiff (P \eif Q)  \text{ e}\\
	& \vdash (\enot\enot P \eand Q) \eiff (P \eand Q)
\end{align*}
e compare as duas.)

Eis outro fato: se duas sentenças $\metav{P}$ e~$\metav{Q}$ são equivalentes e você substitui a mesma letra sentencial em ambas por uma mesma sentença~$\metav{R}$, os resultados também são equivalentes. Por exemplo, se você substituir $A$ em $A \land B$ e $B \land A$ por, digamos, $\enot C$, obterá $\enot C \land B$ e $B \land \enot C$, que são equivalentes. Também podemos registrar isso:

\factoidbox{A equivalência é preservada sob a substituição de letras sentenciais; isto é, se $\metav{P}(A)$ e $\metav{Q}(A)$ contêm ambas a letra sentencial $A$ e são equivalentes, então as sentenças $\metav{P}(\metav{R})$ e $\metav{Q}(\metav{R})$ (obtidas substituindo $A$ por $\metav{R}$ em ambas) também são equivalentes.}

Isso significa que, uma vez demonstrado que duas sentenças são equivalentes (por exemplo, $\enot\enot A$ e $A$, ou $A \land B$ e $B \land A$), sabemos que todas as suas «instâncias» comuns também são equivalentes. Observe que não obtemos isso imediatamente de uma tabela-verdade ou de uma demonstração por dedução natural. Por exemplo, uma tabela-verdade que mostra que $\enot\enot A$ e~$A$ são equivalentes \emph{não} mostra também que $\enot\enot(B \eif C)$ e $B \eif C$ são equivalentes: a primeira exige apenas 2 linhas, enquanto a segunda exige 4.

\section{Cadeias de equivalências}

Quando quiser verificar que duas sentenças são equivalentes, você pode, é claro, fazer uma tabela-verdade ou procurar uma demonstração formal. Mas há um método mais simples, baseado no princípio da substituibilidade de equivalentes que acabamos de discutir: munido de um pequeno catálogo de equivalências simples, substitua partes da primeira sentença por partes equivalentes e repita até chegar à segunda sentença.

Esse método de mostrar que sentenças são equivalentes é fundamentado nos dois fatos da seção anterior. O primeiro fato diz que, \emph{se}, por exemplo, $\enot\enot \metav{P}$ e $\metav{P}$ são equivalentes (para qualquer sentença~$\metav{P}$), substituir $\enot\enot\metav{P}$ por $\metav{P}$ em uma sentença resulta em uma sentença equivalente. O segundo fato diz \emph{que} $\enot\enot \metav{P}$ e~$\metav{P}$ são sempre equivalentes, para qualquer sentença~$\metav{P}$. (Uma simples tabela-verdade mostra que $\enot\enot A$ e~$A$ são equivalentes.) Pelo segundo fato, sabemos que, sempre que substituímos $A$ em $\enot\enot A$ e $A$ por alguma sentença~$\metav{P}$, obtemos resultados equivalentes. Em outras palavras, do fato de $\enot\enot A$ e $A$ serem equivalentes e do segundo fato, podemos concluir que, para qualquer sentença~$\metav{P}$, $\enot\enot \metav{P}$ e~$\metav{P}$ são equivalentes.

Vejamos um exemplo. Pelas leis de De Morgan, os seguintes pares de sentenças são equivalentes:
\begin{align*}
	\enot (A \eand B) & \Leftrightarrow (\enot A \eor \enot B)\\
	\enot (A \eor B) & \Leftrightarrow (\enot A \eand \enot B)
\end{align*}
Isso pode ser verificado construindo duas tabelas-verdade ou quatro demonstrações de dedução natural que mostrem:
\begin{align*}
	\enot (A \eand B) & \vdash (\enot A \eor \enot B)\\
	(\enot A \eor \enot B) & \vdash \enot (A \eand B)\\
	\enot (A \eor B) & \vdash (\enot A \eand \enot B)\\
	(\enot A \eand \enot B) & \vdash \enot (A \eor B)
\end{align*}

Pelo segundo fato, \emph{quaisquer} pares de sentenças das formas a seguir são equivalentes:
\begin{align*}
	\enot (\metav{P} \eand \metav{Q}) & \Leftrightarrow (\enot \metav{P} \eor \enot \metav{Q})\\
	\enot (\metav{P} \eor \metav{Q}) & \Leftrightarrow (\enot \metav{P} \eand \enot \metav{Q})
\end{align*}
Agora considere a sentença $\enot(R \eor (S \eand T))$. Encontraremos uma sentença equivalente na qual todos os sinais $\enot$ se aplicam diretamente a letras sentenciais. No primeiro passo, consideramos isso como uma sentença da forma $\enot(\metav{P} \eor \metav{Q})$; então $\metav{P}$ é a sentença $R$ e $\metav{Q}$ é $(S \eand T)$. Como $\enot(\metav{P} \eor \metav{Q})$ é equivalente a $(\enot \metav{P} \eand \enot \metav{Q})$ (pela segunda lei de De Morgan), podemos substituir a sentença inteira por $(\enot \metav{P} \eand \enot \metav{Q})$. Neste caso, em que $\metav{P}$ é $R$ e $\metav{Q}$ é $(S \eand T)$, obtemos $(\enot R \eand \enot (S \eand T))$. Essa nova sentença contém como parte a sentença $\enot (S \eand T)$. Ela tem a forma $\enot(\metav{P} \eand \metav{Q})$, mas agora $\metav{P}$ é a letra sentencial $S$ e $\metav{Q}$ é $T$. Pela lei de De Morgan (a primeira desta vez), isso é equivalente a $(\enot\metav{P} \eor \enot\metav{Q})$, ou, neste caso específico, a $(\enot S \eor \enot T)$. Portanto, podemos substituir a parte $\enot (S \eand T)$ por $(\enot S \eor \enot T)$. O resultado é a sentença $(\enot R \eand (\enot S \eor \enot T))$, na qual todos os símbolos $\enot$ se aplicam diretamente a letras sentenciais. Nós «empurramos» as negações para dentro tanto quanto possível. Podemos registrar uma cadeia de equivalências assim listando os passos individuais e registrando, como fazemos na dedução natural, qual equivalência básica usamos em cada caso:
\begin{align*}
	& \fbox{\enot(\fbox{$R$} \eor \fbox{$(S \eand T)$})} \\
	& \fbox{(\enot(\fbox{$R$} \eand \enot\fbox{$(S \eand T)$})} && \text{DeM}\\
	& (\enot(R \eand \fbox{$\enot(\fbox{$S$} \eand \fbox{$T$})$}) \\
	& (\enot(R \eand \fbox{$(\enot \fbox{$S$} \eor \enot \fbox{$T$})$}) && \text{DeM}
\end{align*}
Destacamos a sentença substituída e aquelas que correspondem a $\metav{P}$ e~$\metav{Q}$ nas leis de De Morgan para facilitar a leitura, mas isso não é necessário e não continuaremos a fazê-lo.

Em \cref{tab:equivalences}, apresentamos uma lista de equivalências básicas que você pode usar nessas cadeias de equivalências. Os rótulos abreviam os nomes tradicionais das respectivas leis lógicas: dupla negação (DN), De Morgan (DeM), comutatividade (Comm), distributividade (Dist), associatividade (Assoc), idempotência (Id) e absorção (Abs).

\begin{table}
	\begin{align*}
\lnot\lnot \metav{P} & \Leftrightarrow \metav{P} && \text{(DN)}\\[2ex]
(\metav{P} \eif \metav{Q}) & \Leftrightarrow (\lnot \metav{P} \lor \metav{Q})
&& \text{(Cond)}\\
\lnot(\metav{P} \eif \metav{Q}) & \Leftrightarrow (\metav{P} \land \lnot\metav{Q}) \\[2ex]
(\metav{P} \eiff \metav{Q}) & \Leftrightarrow ((\metav{P} \eif \metav{Q}) \land  (\metav{Q} \eif \metav{P}))
&& \text{(Bicond)}\\[2ex]
\lnot(\metav{P} \land \metav{Q}) & \Leftrightarrow (\lnot\metav{P} \lor \lnot\metav{Q})
&& \text{(DeM)}\\
\lnot(\metav{P} \lor \metav{Q}) & \Leftrightarrow (\lnot\metav{P} \land \lnot\metav{Q}) \\[2ex]
(\metav{P} \lor \metav{Q}) & \Leftrightarrow (\metav{Q} \lor \metav{P}) & &\text{(Comm)}\\
(\metav{P} \land \metav{Q}) & \Leftrightarrow (\metav{Q} \land \metav{P})\\[2ex]
(\metav{P} \land (\metav{Q} \lor \metav{R})) & \Leftrightarrow ((\metav{P} \land \metav{Q}) \lor (\metav{P} \land \metav{R}))
&& \text{(Dist)}\\
(\metav{P} \lor (\metav{Q} \land \metav{R})) & \Leftrightarrow ((\metav{P} \lor \metav{Q}) \land (\metav{P} \lor \metav{R}))\\
(\metav{P} \lor (\metav{Q} \lor \metav{R})) & \Leftrightarrow ((\metav{P} \lor \metav{Q}) \lor \metav{R}) & &\text{(Assoc)}\\
(\metav{P} \land (\metav{Q} \land \metav{R})) & \Leftrightarrow ((\metav{P} \land \metav{Q}) \land \metav{R})\\[2ex]
(\metav{P} \lor \metav{P}) & \Leftrightarrow \metav{P} && \text{(Id)}\\
(\metav{P} \land \metav{P}) & \Leftrightarrow \metav{P}\\[2ex]
(\metav{P} \land (\metav{P} \lor \metav{Q})) & \Leftrightarrow \metav{P}&& \text{(Abs)}\\
(\metav{P} \lor (\metav{P} \land \metav{Q})) & \Leftrightarrow \metav{P}\\[2ex]
(\metav{P} \land (\metav{Q} \lor \lnot\metav{Q})) & \Leftrightarrow \metav{P}  && \text{(Simp)}\\
(\metav{P} \lor (\metav{Q} \land \lnot\metav{Q})) & \Leftrightarrow \metav{P}\\[2ex]
(\metav{P} \lor (\metav{Q} \lor \lnot\metav{Q})) & \Leftrightarrow (\metav{Q} \lor \lnot\metav{Q}) \\
(\metav{P} \land (\metav{Q} \land \lnot\metav{Q})) & \Leftrightarrow (\metav{Q} \land \lnot\metav{Q})
\end{align*}
\caption{Equivalências básicas}
\label{tab:equivalences}
\end{table}

\section{Encontrando formas normais equivalentes}

Em \cref{c:NormalForms}, mostramos que toda sentença da LVF é equivalente a uma sentença na forma normal disjuntiva (FND) e a uma na forma normal conjuntiva (FNC). Fizemos isso apresentando um método para construir uma sentença na FND ou na FNC equivalente à sentença original: primeiro construímos uma tabela-verdade e depois «lemos» nela uma sentença na FND ou na FNC. Esse método tem duas desvantagens. A primeira é que as sentenças resultantes na FND ou na FNC nem sempre são as mais curtas. A segunda é que o próprio método se torna difícil de aplicar quando a sentença inicial contém mais do que algumas letras sentenciais (pois a tabela-verdade de uma sentença com $n$ letras sentenciais tem $2^n$ linhas).

Podemos usar cadeias de equivalências como método alternativo: para encontrar uma sentença na FND, podemos aplicar sucessivamente equivalências básicas até encontrar uma sentença equivalente que esteja na FND. Lembre-se das condições que uma sentença na FND deve satisfazer:
\begin{compactlist}
	\item[(\textsc{dnf1})] Não ocorrem na sentença conectivos além de negações, conjunções e disjunções;
	\item[(\textsc{dnf2})] Toda ocorrência de negação tem escopo mínimo (isto é, qualquer ocorrência de $\enot$ é imediatamente seguida por uma sentença atômica);
	\item[(\textsc{dnf3})] Nenhuma disjunção ocorre no escopo de uma conjunção.
\end{compactlist}
A condição (\textsc{dnf1}) diz que precisamos remover de uma sentença todos os símbolos $\eif$ e~$\eiff$. É para isso que servem as equivalências básicas (Cond) e (Bicond). Por exemplo, suponha que começamos com a sentença
\begin{align*}
	& \enot(A \eand \enot C) \land (\enot A\eif \enot B).
\intertext{Podemos eliminar $\eif$ usando (Cond). Neste caso, $\metav{P}$ é $\enot A$ e $\metav{Q}$ é $\enot{B}$. Obtemos:}
& \enot(A \eand \enot C) \land (\enot\enot A \eor \enot B) && \text{Cond}
\intertext{A dupla negação pode ser removida, pois $\enot\enot A$ é equivalente a~$A$:}
& \enot(A \eand \enot C) \land (A \eor \enot B) && \text{DN}
\intertext{Agora a condição (\textsc{dnf1}) está satisfeita: nossa sentença contém apenas $\enot$, $\eand$ e $\eor$. A condição (\textsc{dnf2}) diz que precisamos encontrar uma maneira de fazer todos os $\enot$ se aplicarem imediatamente a letras sentenciais. Mas isso não ocorre na primeira conjunção. Para garantir que (\textsc{dnf2}) seja satisfeita, usamos as leis de De Morgan e a lei da dupla negação (DN) tantas vezes quantas forem necessárias.}
& (\enot A \eor \enot\enot C) \land (A \eor \enot B) && \text{DeM}\\
& (\enot A \eor C) \land (A \eor \enot B) && \text{DN}
\intertext{A sentença resultante está agora na FNC: é uma conjunção de disjunções de letras sentenciais e letras sentenciais negadas. Mas queremos uma sentença na FND, isto é, uma sentença em que (\textsc{dnf3}) seja satisfeita: nenhum $\eor$ ocorre no escopo de um $\eand$. Usamos as leis distributivas (Dist) para garantir isso. A última sentença tem a forma $\metav{P} \land (\metav{Q} \lor \metav{R})$, em que $\metav{P}$ é $(\enot A \eor C)$, $\metav{Q}$ é $A$ e $\metav{R}$ é $\enot B$. Aplicando (Dist) uma vez, obtemos:}
& ((\enot A \eor C) \eand A)) \eor ((\enot A \eor C) \eand \enot B) && \text{Dist}
\intertext{Isso parece pior, mas, se continuarmos, vai melhorar! As duas disjunções quase parecem permitir a aplicação de (Dist) novamente, exceto pelo fato de que $\eor$ está do lado errado. É para isso que serve a comutatividade (Comm). Vamos aplicá-la a $(\enot A \eor C) \eand A$:}
& (A \eand (\enot A \eor C)) \eor ((\enot A \eor C) \eand \enot B) && \text{Comm}
\intertext{Podemos aplicar (Dist) novamente à parte resultante, $A \eand (\enot A \eor C)$:}
& ((A \eand \enot A) \eor (A \eand C)) \eor ((\enot A \eor C) \eand \enot B) && \text{Dist}
\intertext{Agora, na metade esquerda, nenhum $\eor$ está no escopo de um $\eand$. Aplicaremos os mesmos princípios à metade direita:}
& ((A \eand \enot A) \eor (A \eand C)) \eor (\enot B \eand (\enot A \eor C)) && \text{Comm}\\
& ((A \eand \enot A) \eor (A \eand C)) \eor ((\enot B \eand \enot A) \eor (\enot B \eand C)) && \text{Dist}
\intertext{Nossa sentença está agora na FND! Mas podemos simplificá-la um pouco: $(A \eand \enot A)$ é uma contradição na LVF, isto é, é sempre falsa. E, se combinamos algo que é sempre falso usando $\eor$ com uma sentença~$\metav{P}$, obtemos algo equivalente simplesmente a~$\metav{P}$. Esta é a segunda das regras de simplificação (Simp).}
& ((A \eand C) \eor (A \eand \enot A)) \eor ((\enot B \eand \enot A) \eor (\enot B \eand C)) && \text{Comm}\\
& (A \eand C) \eor ((\enot B \eand \enot A) \eor (\enot B \eand C)) && \text{Simp}
\end{align*}
O resultado final ainda está na FND, mas é um pouco mais simples. Também é muito mais simples que a FND que teríamos obtido pelo método de \cref{c:NormalForms}. Na verdade, a sentença com que começamos poderia ser a $\metav{S}$ de \cref{s:DNFTruthTable}: ela tem exatamente a tabela-verdade usada como exemplo ali. A FND que encontramos ali (\ifHTMLtarget na seção~\cref{s:DNFTruthTable}\else na p.~\pageref{longDNF}\fi) era, com todos os parênteses necessários:
$$((((A \eand B) \eand C) \eor ((A \eand \enot B) \eand C)) \eor ((\enot A \eand \enot B) \eand C)) \eor ((\enot A \eand \enot B) \eand \enot C)$$

\practiceproblems
\problempart
\label{pr.DNF2}
Considere as seguintes sentenças:
\begin{compactlist}
	\item $(A \eif \enot B)$
	\item $\enot (A \eiff B)$
	\item $(\enot A \eor \enot (A \eand B))$
	\item $(\enot (A \eif B ) \eand (A \eif C))$
	\item $(\enot (A \eor B) \eiff ((\enot C \eand \enot A) \eif \enot B))$
	\item $((\enot (A \eand \enot B) \eif C) \eand \enot (A \eand D))$
\end{compactlist}
Para cada sentença, encontre uma sentença equivalente na FND e outra na FNC fornecendo uma cadeia de equivalências. Use (Id), (Abs) e (Simp) para simplificar suas sentenças tanto quanto possível.

\chapter{Correção}\label{ch:Soundness}

Neste capítulo, relacionamos a semântica da LVF ao seu \emph{sistema de provas} de dedução natural (definido em \cref{ch.NDTFL}). Demonstraremos que o sistema formal de provas é seguro: só é possível demonstrar, a partir de premissas, sentenças que de fato decorrem delas.
Intuitivamente, um sistema formal de provas é correto \emph{\ifeff} não permite demonstrar argumentos inválidos. Essa é obviamente uma propriedade muito desejável. Ela nos diz que nosso sistema de provas nunca nos levará ao erro. De fato, se nosso sistema não fosse correto, não poderíamos confiar em nossas demonstrações. O objetivo deste capítulo é provar que nosso sistema de provas é correto.

Tornemos a ideia mais precisa. Abreviaremos uma lista de sentenças usando a letra grega $\Gamma$ («gama»). Um sistema formal de provas é \define{correto} (relativamente a uma semântica dada) \emph{\ifeff} sempre que houver uma demonstração formal de $\metav{C}$ a partir de suposições entre $\Gamma$, então $\Gamma$ implica genuinamente $\metav{C}$ (dada essa semântica). Em outras palavras, para provar que o sistema de provas da LVF é correto, precisamos provar o seguinte

\begin{factoidboxe}\textbf{Teorema da correção.} Para quaisquer sentenças $\Gamma$ e $\metav{C}$: se $\Gamma\proves\metav{C}$, então $\Gamma \entails\metav{C}$
\end{factoidboxe}

Para prová-lo, verificaremos individualmente cada regra do sistema de provas da LVF. Queremos mostrar que nenhuma aplicação dessas regras nos leva ao erro. Como uma demonstração envolve apenas aplicações repetidas dessas regras, isso mostrará que nenhuma demonstração nos leva ao erro. Ou, pelo menos, essa é a ideia geral.

Para começar, precisamos tornar mais precisa a ideia de «levar ao erro». Diremos que uma linha de uma demonstração é \define{brilhante} \emph{\ifeff} as suposições das quais essa linha depende implicam a sentença escrita na linha.\footnote{A palavra «brilhante» não é padrão entre os lógicos.} Para ilustrar a ideia, considere:
	\begin{fitchproof}
		\hypo{fgh}{F\eif(G\eand H)}\PR
		\open
			\hypo{f}{F}\AS
			\have{gh}{G \eand H}\ce{fgh,f}
			\have{g}{G}\ae{gh}
		\close
		\have{fg}{F \eif G}\ci{f-g}
	\end{fitchproof}\noindent\noindent
A linha $1$ é brilhante \emph{\ifeff} $F \eif (G \eand H) \entails F \eif (G \eand H)$. Você deve estar facilmente convencido de que a linha $1$ é, de fato, brilhante! Do mesmo modo, a linha $4$ é brilhante \emph{\ifeff} $F \eif (G \eand H), F \entails G$. Novamente, é fácil verificar que a linha~$4$ é brilhante. O mesmo vale para todas as linhas desta demonstração na LVF. Queremos mostrar que isso não é coincidência. Isto é, queremos provar:
	\begin{factoidboxe}\textbf{Lema do brilho.}
		Toda linha de toda demonstração na LVF é brilhante.
	\end{factoidboxe}\noindent
Então saberemos que nunca nos desviamos, em nenhuma linha de uma demonstração. De fato, dado o lema do brilho, será fácil provar o teorema da correção:

\emph{Demonstração.} Suponha que $\Gamma \proves \metav{C}$. Então existe uma demonstração na LVF, com $\metav{C}$ aparecendo em sua última linha, cujas únicas suposições não descarregadas estão entre $\Gamma$. O lema do brilho diz que toda linha de toda demonstração na LVF é brilhante. Logo, essa última linha é brilhante, isto é, $\Gamma \entails \metav{C}$. QED

Resta demonstrar o lema do brilho.

Para isso, observamos que toda linha de qualquer demonstração na LVF é uma premissa ou uma suposição, ou é obtida aplicando alguma regra.

As premissas são automaticamente brilhantes: se $\metav{A}$ é uma premissa, então está entre as sentenças de~$\Gamma$, e $\Gamma \entails \metav{A}$ trivialmente. As suposições também são brilhantes, pois qualquer suposição $\metav{A}$ depende de si mesma, e $\metav{A} \entails \metav{A}$. Portanto, o que queremos mostrar é que nenhuma aplicação de uma regra do sistema de provas da LVF nos levará ao erro. Mais precisamente, diremos que uma regra de inferência é \define{correta quanto \`as regras} \emph{\ifeff} para todas as demonstrações na LVF, se obtemos uma linha numa demonstração aplicando essa regra e, quando toda linha anterior da demonstração é brilhante, nossa nova linha também é brilhante. Precisamos mostrar que \emph{toda} regra do sistema de provas da LVF é correta quanto às regras.











Faremos isso a seguir. Mas, tendo demonstrado a correção quanto às regras de todas as regras, o lema do brilho seguirá imediatamente:

\emph{Demonstração.} Comece com a linha~$1$ de qualquer demonstração na LVF. Ela deve ser uma premissa ou uma suposição, e ambas são brilhantes, como vimos acima. Tome a linha seguinte, $2$. Se ela é uma premissa ou suposição, é brilhante. Caso contrário, é obtida da linha~$1$ usando uma inferência correta quanto às regras. Como a linha~$1$ é brilhante, a linha~$2$ também é brilhante.




Tome a linha seguinte, $3$. Se ela é uma premissa ou suposição, é brilhante. Caso contrário, é obtida de uma linha anterior usando uma inferência correta quanto às regras, e estabelecemos que todas as linhas anteriores são brilhantes. Assim, a linha~$3$ também é brilhante.



E assim por diante. Em geral, a sentença escrita na linha $n$ deve ser uma premissa ou suposição (que é brilhante) ou ser obtida usando uma regra formal de inferência correta quanto às regras. Como toda linha anterior é brilhante, a própria linha~$n$ é brilhante. Podemos percorrer esse raciocínio, para qualquer demonstração na LVF, começando com a linha~$1$ e continuando até a última linha, e concluir que toda linha de toda demonstração na LVF é brilhante.





QED

Resta mostrar que toda regra é correta quanto às regras. Isso não é difícil, mas é demorado, pois precisamos verificar cada regra individualmente, e o sistema de provas da LVF tem muitas regras! Para acelerar um pouco o processo, introduziremos uma abreviação conveniente:



$\Delta_i$ («delta») abreviará as suposições (se houver) das quais a linha~$i$ depende em nossa demonstração na LVF (o contexto indicará qual demonstração na LVF temos em mente). Isso inclui todas as premissas de nossa demonstração e todas as suposições de subdemonstrações que ainda estão abertas na linha~$i$.



Primeiro, registremos nossa observação anterior sobre premissas e suposições.

\begin{factoidboxe}
	Premissas e suposições em demonstrações na LVF são brilhantes.
\end{factoidboxe}

Se $\metav{A}$ é uma premissa na linha~$n$, então está entre $\Delta_n$, pois isso inclui todas as premissas da demonstração. Se é introduzida como suposição de uma subdemonstração na linha~$n$, então tudo na subdemonstração (incluindo a linha~$n$, isto é, a própria $\metav{A}$) depende de~$\metav{A}$ e, portanto, $\metav{A}$ está entre $\Delta_n$. Em qualquer caso, $\Delta_n \entails \metav{A}$.






Agora prossigamos para mostrar que todas as regras de inferência são corretas quanto às regras.

\begin{factoidboxe}$\eand$I é correta quanto às regras.
\end{factoidboxe}

\emph{Demonstração.} Considere qualquer aplicação de $\eand$I em qualquer demonstração na LVF, isto é, algo como:
\begin{fitchproof}
	\have[i]{a}{\metav{A}}
	\have[j]{b}{\metav{B}}
	\have[n]{c}{\metav{A}\eand\metav{B}} \ai{a, b}
\end{fitchproof}\noindent
Para mostrar que $\eand$I é correta quanto às regras, supomos que toda linha anterior à linha $n$ é brilhante e pretendemos mostrar que a linha $n$ é brilhante, isto é, que $\Delta_n \entails \metav{A} \eand \metav{B}$.

Seja então $v$ uma valoração qualquer que torna verdadeiras todas as sentenças de $\Delta_n$.

Primeiro mostramos que $v$ torna $\metav{A}$ verdadeira. Para isso, observe que todas as sentenças de $\Delta_i$ estão entre as de $\Delta_n$. Pela hipótese, a linha $i$ é brilhante. Assim, qualquer valoração que torne verdadeiras todas as sentenças de $\Delta_i$ torna $\metav{A}$ verdadeira. Como $v$ torna verdadeiras todas as sentenças de $\Delta_i$, também torna $\metav{A}$ verdadeira.

De modo semelhante, vemos que $v$ torna $\metav{B}$ verdadeira.

Logo, $v$ torna $\metav{A}$ verdadeira e $v$ torna $\metav{B}$ verdadeira. Consequentemente, $v$ torna $\metav{A}\eand\metav{B}$ verdadeira. Portanto, qualquer valoração que torne verdadeiras todas as sentenças de $\Delta_n$ também torna $\metav{A} \eand \metav{B}$ verdadeira. Isto é: a linha $n$ é brilhante. QED


Todos os lemas restantes que estabelecem a correção quanto às regras terão, essencialmente, a mesma estrutura que este.

\begin{factoidboxe}$\eand$E é correta quanto às regras.
\end{factoidboxe}

\emph{Demonstração.}
	Suponha que toda linha anterior à linha $n$ em alguma demonstração na LVF é brilhante e que $\eand$E é usada na linha $n$. A situação é:
		   \begin{fitchproof}
			   \have[i]{ab}{\metav{A}\eand\metav{B}}
			   \have[n]{a}{\metav{A}} \ae{ab}
		   \end{fitchproof}\noindent
(ou talvez com $\metav{B}$ na linha $n$ em vez disso; mas um raciocínio semelhante se aplicará nesse caso). Seja $v$ uma valoração qualquer que torne verdadeiras todas as sentenças de $\Delta_n$. Observe que todas as sentenças de $\Delta_i$ estão entre as de $\Delta_n$. Pela hipótese, a linha $i$ é brilhante. Assim, qualquer valoração que torne verdadeiras todas as sentenças de $\Delta_i$ torna $\metav{A}\eand\metav{B}$ verdadeira. Logo, $v$ torna $\metav{A}\eand\metav{B}$ verdadeira e, portanto, torna $\metav{A}$ verdadeira. Assim, $\Delta_n \entails \metav{A}$. QED


\begin{factoidboxe}$\eor$I é correta quanto às regras.
\end{factoidboxe}

Deixamos isso como exercício.

\begin{factoidboxe}$\eor$E é correta quanto às regras.
\end{factoidboxe}

\emph{Demonstração.}
	Suponha que toda linha anterior à linha $n$ em alguma demonstração na LVF é brilhante e que $\eor$E é usada na linha $n$. A situação é:
   \begin{fitchproof}
	   \have[m]{aob}{\metav{A}\eor\metav{B}}
	   \open
		   \hypo[i]{a}{\metav{A}}\AS %\by{want \metav{C}}{}
		   \have[j]{c1}{\metav{C}}
	   \close
	   \open
		   \hypo[k]{b}{\metav{B}}\AS %\by{want \metav{C}}{}
		   \have[l]{c2}{\metav{C}}
	   \close
	   \have[n]{ab}{\metav{C}}\oe{aob, a-c1,b-c2}
   \end{fitchproof}\noindent
Seja $v$ uma valoração qualquer que torne verdadeiras todas as sentenças de $\Delta_n$. Observe que todas as sentenças de $\Delta_m$ estão entre as de $\Delta_n$. Pela hipótese, a linha $m$ é brilhante. Assim, qualquer valoração que torne verdadeiras todas as sentenças de $\Delta_n$ torna $\metav{A} \eor \metav{B}$ verdadeira. Em particular, $v$ torna $\metav{A} \eor \metav{B}$ verdadeira e, portanto, ou $v$ torna $\metav{A}$ verdadeira ou $v$ torna $\metav{B}$ verdadeira. Agora analisamos esses dois casos:
   \begin{numberlist}
	   \item \emph{$v$ torna $\metav{A}$ verdadeira.} Todas as sentenças de $\Delta_i$ estão entre as de $\Delta_n$, com a possível exceção de $\metav{A}$. Como $v$ torna verdadeiras todas as sentenças de $\Delta_n$ e também torna $\metav{A}$ verdadeira, $v$ torna verdadeiras todas as sentenças de $\Delta_i$. Pela hipótese, a linha $j$ é brilhante; portanto, $\Delta_j \entails \metav{C}$. Mas as sentenças de $\Delta_i$ são precisamente as sentenças de $\Delta_j$, logo $\Delta_i \entails \metav{C}$. Assim, qualquer valoração que torne verdadeiras todas as sentenças de $\Delta_i$ torna $\metav{C}$ verdadeira. Mas $v$ é justamente uma valoração desse tipo. Portanto, $v$ torna $\metav{C}$ verdadeira.
	   \item \emph{$v$ torna $\metav{B}$ verdadeira.} Raciocinando exatamente da mesma maneira, considerando as linhas $k$ e~$l$, $v$ torna $\metav{C}$ verdadeira.
	   \end{numberlist}
De qualquer modo, $v$ torna $\metav{C}$ verdadeira. Portanto, $\Delta_n \entails \metav{C}$.
QED


\begin{factoidboxe}
	$\enot$E é correta quanto às regras.
\end{factoidboxe}

\emph{Demonstração.}
	Suponha que toda linha anterior à linha $n$ em alguma demonstração na LVF é brilhante e que $\enot$E é usada na linha $n$. A situação é:
\begin{fitchproof}
   \have[i]{i}{\metav{A}}
   \have[j]{j}{\enot\metav{A}}
   \have[n]{nb}{\ered}\ri{i, j}
\end{fitchproof}\noindent
Observe que todas as sentenças de $\Delta_i$ e todas as de $\Delta_j$ estão entre as de $\Delta_n$. Pela hipótese, as linhas $i$ e $j$ são brilhantes. Assim, qualquer valoração que torne verdadeiras todas as sentenças de $\Delta_n$ teria de tornar $\metav{A}$ e $\enot\metav{A}$ verdadeiras. Mas nenhuma valoração pode fazer isso. Portanto, nenhuma valoração torna verdadeiras todas as sentenças de $\Delta_n$. Logo, $\Delta_n \entails \ered$, vacuamente.
QED

\begin{factoidboxe}
	X é correta quanto às regras.
\end{factoidboxe}

Deixamos isso como exercício.

\begin{factoidboxe}
	$\enot$I é correta quanto às regras.
\end{factoidboxe}

\emph{Demonstração.}
	Suponha que toda linha anterior à linha $n$ em alguma demonstração na LVF é brilhante e que $\enot$I é usada na linha $n$. A situação é:
\begin{fitchproof}
   \open
	   \hypo[i]{a}{\metav{A}}\AS
	   \have[j]{b}{\ered}
   \close
   \have[n]{na}{\enot\metav{A}}\ni{a-b}
\end{fitchproof}\noindent
Seja $v$ uma valoração qualquer que torne verdadeiras todas as sentenças de $\Delta_n$. Observe que todas as sentenças de $\Delta_n$ estão entre as de $\Delta_i$, com a possível exceção da própria $\metav{A}$. Pela hipótese, a linha $j$ é brilhante. Mas nenhuma valoração pode tornar $\ered$ verdadeira, portanto nenhuma valoração pode tornar verdadeiras todas as sentenças de $\Delta_j$. Como as sentenças de $\Delta_i$ são precisamente as de $\Delta_j$, nenhuma valoração pode tornar verdadeiras todas as sentenças de $\Delta_i$. Como $v$ torna verdadeiras todas as sentenças de $\Delta_n$, ela deve tornar $\metav{A}$ falsa e, portanto, tornar $\enot \metav{A}$ verdadeira. Assim, $\Delta_n \entails \enot \metav{A}$.
QED


\begin{factoidboxe}\label{lem:LastRuleSound} IP, $\eif$I, $\eif$E, $\eiff$I e $\eiff$E são todas corretas quanto às regras.
\end{factoidboxe}

Deixamos essas demonstrações como exercícios.

Isso estabelece que todas as regras básicas de nosso sistema de provas são corretas quanto às regras. Por fim, mostramos:

\begin{factoidboxe}Todas as regras derivadas de nosso sistema de provas são corretas quanto às regras.
\end{factoidboxe}

\emph{Demonstração.}
	Suponha que usamos uma regra derivada para obter alguma sentença, $\metav{A}$, na linha $n$ de uma demonstração na LVF e que toda linha anterior é brilhante. Cada uso de uma regra derivada pode ser substituído (ao custo de uma longa explicação) por vários usos de regras básicas. Isto é, poderíamos ter usado regras básicas para escrever $\metav{A}$ em alguma linha $n+k$, sem introduzir suposições adicionais. Assim, aplicando várias vezes nossos resultados individuais de que todas as regras básicas são corretas quanto às regras (na verdade, $k+1$ vezes), podemos ver que a linha $n+k$ é brilhante. Portanto, a regra derivada é correta quanto às regras.
QED


E é isso! Mostramos que toda regra — básica ou não — é correta quanto às regras, que é tudo o que precisávamos para estabelecer o lema do brilho e, consequentemente, o teorema da correção.

Pode ser útil encerrar este capítulo repetindo minha explicação informal do que fizemos. Uma demonstração formal é apenas uma sequência — de comprimento arbitrário — de aplicações de regras. Mostramos que nenhuma aplicação de nenhuma regra nos levará ao erro. Segue-se que nenhuma demonstração formal nos levará ao erro.




Isto é: nosso sistema de provas é correto.

\practiceproblems

\problempart
\label{pr.Soundness}
Complete os lemas deixados como exercícios neste capítulo. Isto é, mostre que são corretas quanto às regras:
	\begin{compactlist}
		\item $\eor$I. (\emph{Dica}: é semelhante ao caso de $\eand$E.)
		\item X. (\emph{Dica}: é semelhante ao caso de $\enot$E.)
		\item $\eif$I. (\emph{Dica}: é semelhante a $\eor$E.)
		\item $\eif$E.
		\item IP. (\emph{Dica}: é semelhante ao caso de $\enot$I.)
	\end{compactlist}

